\BourbakiExerciseGroup

\begin{enumerate}
\def\labelenumi{\arabic{enumi})}
\tightlist
\item
若 \(x\) 与 \(y\) 是实数，我们有
\end{enumerate}

\[
[ x + y ] = [ x ] + [ y ] + \varepsilon \quad \text { where } \varepsilon = 0 \quad \text { or } \quad \varepsilon = 1,
\]

\[
[ x - y ] = [ x ] - [ y ] - \varepsilon \quad \text { where } \varepsilon = 0 \quad \text { or } \quad \varepsilon = 1,
\]

\[
[ x ] + [ x + y ] + [ y ] \leqslant [ 2 x ] + [ 2 y ],
\]

\[
[ x ] + \left[ x + \frac {1}{n} \right] + \left[ x + \frac {2}{n} \right] + \dots + \left[ x + \frac {n - 1}{n} \right] = [ n x ],
\]

\[
\left[ \frac {[ n x ]}{n} \right] = [ x ]
\]

对任何整数 \(n > 0\) 成立。

\begin{enumerate}
\def\labelenumi{\arabic{enumi})}
\setcounter{enumi}{1}
\item
设 \((\varepsilon_{n})\) 是趋于 0 的有限数 \(>0\) 的一个严格递减序列。对每个 \(x \in \mathbb{R}\) ，令 \(r_{n}(x)\) 为 \(\varepsilon_{n}\) 的那个倍数，它是 \(x\) 的精确到 \(\varepsilon_{n}\) 的不足近似。则序列 \((r_{n}(x))\) 对一切 \(x \in \mathbb{R}\) 递增，当且仅当 \(\varepsilon_{n}\) 是 \(\varepsilon_{n+1}\) 的整数倍（对每个 \(n\) ）。
\item
设 \(a\) 是一个整数 \(> 1\) 。实数 \(x\) 是有理数，当且仅当它以 \(a\) 为底的展开式，比如说 \(x = p_0 + \sum_{n=0}^{\infty} u_n a^{-n}\) ，是周期的，也就是说，当且仅当存在两个整数 \(n_0\) 与 \(r > 0\) ，使得 \(u_{n+r} = u_n\) 对一切 \(n \geqslant n_0\) 成立。
\item
设 \((u_{n})\) 是数 \(>0\) 的一个可和序列（在 \(\mathbf{R}\) 中），满足下列条件：\(u_{n+1} \leqslant u_{n}\) 且 \(u_{n} \leqslant \sum_{k=1}^{\infty} u_{n+k}\) 对一切 \(n \geqslant 0\) 成立。试证：对每个数 \(a\) （满足 \(0 < a \leqslant s = \sum_{n=0}^{\infty} u_{n}\) ），存在 \(\mathbf{N}\) 的一个子集 I，使得 \(a = \sum_{n \in \mathbf{I}} u_{n}\) 。若 \(u_{n+1} \leqslant u_{n} \leqslant 2u_{n+1}\) 对一切 \(n\) 成立，这些条件特别地被满足。考察二进展开的情形。
\item
对每个实数 \(x \in ]0, I]\) ，存在唯一的无穷递增序列 \((q_n)\) （整数 \(> 0\) 的），使得 \(x = \sum_{n=1}^{\infty} I / (q_1 q_2 \ldots q_n)\) 。数 \(x\) 是有理数，当且仅当 \(q_{n+1} = q_n\) 从 \(n\) 的某个值起成立。
\item
对每个实数 \(x > \mathbf{I}\) ，存在唯一的无穷序列 \((t_n)\) （整数 \(\geqslant \mathbf{I}\) 的），使得 \(t_{n+1} \geqslant t_n^2\) 对一切 \(n \geqslant \mathbf{I}\) 成立，且使得 \(x = \prod_{n=1}^{\infty} \left( \mathbf{I} + \frac{\mathbf{I}}{t_n} \right)\) （用 §7 习题 19）。数 \(x\) 是有理数，当且仅当 \(t_{n+1} = t_n^2\) 从 \(n\) 的某个值起成立。{[}为证条件是必要的，试证：若令 \(x_n = \prod_{k=n+1}^{\infty} \left( \mathbf{I} + \frac{\mathbf{I}}{t_k} \right)\) ，则分式 \(\frac{x_n}{x_n - \mathbf{I}}\) （取既约形式）的分母构成一个递减序列{]}。
\end{enumerate}

¶ * 7) 设 \((\varepsilon_n)_{n \geqslant 0}\) 是一个每项为 1 或 --- 1 的数的序列。试证：实数

\[
x _ {n} = \varepsilon_ {0} \sqrt {2 + \varepsilon_ {1} \sqrt {2 + \varepsilon_ {2} \sqrt {2 + \cdots + \varepsilon_ {n} \sqrt {2}}}}
\]

有定义且等于

\[
2 \sin \left(\frac {\pi}{4} \sum_ {k = 0} ^ {n} \frac {\varepsilon_ {0} \varepsilon_ {1} \dots \varepsilon_ {k}}{2 ^ {k}}\right).
\]

由此推出：\(\lim_{n\to \infty}x_n\) 存在，且对每个实数 \(x\) （满足 \(-2\leqslant x\leqslant 2\) ），存在一个序列 \((\varepsilon_n)\) ，使得 \(x\) 等于相应序列 \((x_{n})\) 的极限。对哪些 \(x\) 的值序列 \((\varepsilon_n)\) 是唯一的？对哪些 \(x\) 的值它是周期的？（见习题 3）。*

\begin{enumerate}
\def\labelenumi{\arabic{enumi})}
\setcounter{enumi}{7}
\item
试证：不存在非常数映射 \(\psi\) （区间 \(\mathbf{I} \subset \mathbf{R}\) 到自然数序列全体所成之集 \(\mathbf{N}^{\mathbb{N}}\) 中的），它对每个 \(x \in \mathbf{I}\) 具有下列性质：对每个整数 \(n \geqslant 0\) ，存在一个邻域 \(\mathbf{V}\) （ \(x\) 在 \(\mathbf{I}\) 中的），使得对每个 \(y \in \mathbf{V}\) ，前 \(n\) 项（序列 \(\psi(y)\) 的）与前 \(n\) 项（序列 \(\psi(x)\) 的）相同（注意：这将蕴涵 \(\psi\) 的连续性，若每个因子 \(\mathbf{N}\) 带有离散拓扑而 \(\mathbf{N}^{\mathbb{N}}\) 带有积拓扑）。
\item
试证：Cantor 三分集 K（§ 2，第 5 节）是所有这样的实数 \(x \in [0, 1]\) 所成之集，它们的三进展开（相应地，若 \(x\) 是 K 的一个邻接区间的端点，则为非正常三进展开）
\end{enumerate}

\[
x = \sum_ {n = 1} ^ {\infty} u _ {n} 3 ^ {- n}
\]

满足 \(u_{n} \neq \mathrm{I}\) （对每个 \(n\) ，因而 \(u_{n} = 0\) 或 \(u_{n} = 2\) ）。若 \(v_{n} = \frac{\mathrm{I}}{2} u_{n}\) 对每个 \(n\) 成立，且 \(f(x) = \sum_{n=1}^{\infty} v_{n} 2^{-n}\) ，试证：\(f\) 是 \(\mathbf{K}\) 到 \([\mathrm{o}, \mathrm{I}]\) 上的一个满射连续映射；由此推出 \(\mathbf{K}\) 具有连续统的势。

¶ 10) 设 \((d_n)\) 是一个基序列，令 \(a_n = d_n / d_{n-1}\) （ \(n \geqslant 1\) ），设 \((n_i)\) 是整数的一个严格递增序列，并设 \((b_i)\) 是整数的一个序列，满足 \(0 < b_i < a_{n_i} - 1\) （对每个 \(i\) ）。设 \(G\) 为所有这样的 \(x \in [0, 1]\) 所成之集：若 \(\sum_{n=1}^{\infty} u_n / d_n\) 是 \(x\) 的展开式（或非正常展开式，若它存在），则 \(u_{n_i} \neq b_i\) 对一切 \(i \in N\) 成立。试证：\(G\) 是一个全不连通完备集；并且若 \(l\) 是邻接于 \(G\) 且含于 \([0, 1]\) 的诸区间长度之和，则

\[
\mathrm{I} - l = \prod_ {i = 0} ^ {\infty} \left(\mathrm{I} - \frac {\mathrm{I}}{a _ {n _ {i}}}\right).
\]

¶ 11) 设 X 是一个豪斯多夫拓扑空间。假设对每个各项为 o 或 1 的有限序列 s，存在 X 的一个非空子集 \(\mathrm{A}(s)\) ，满足下列条件：(i) 若 s 是 n 项（每项为 o 或 i）的序列，而 \(s'\) , \(s''\) 是 \(n + 1\) 项（每项为 o 或 i）且前 n 项与 s 的相同的序列，则 \(\mathrm{A}(s) = \mathrm{A}(s') \cup \mathrm{A}(s'')\) ；又若 \(s_0\) 是空序列，则 \(\mathrm{A}(s_0) = \mathrm{X}\) 。

\begin{enumerate}
\def\labelenumi{(\roman{enumi})}
\setcounter{enumi}{1}
\tightlist
\item
对每个各项为 0 或 1 的无穷序列 \((u_n)_{n\geqslant 0}\) ，若用 \(s_n\) 表示有限序列 \((u_k)_{0\leqslant k\leqslant n}\) ，则由诸 \(\mathbf{A}(s_n)\) 形成的滤子基收敛于 \(\mathbf{X}\) 的一点。
\end{enumerate}

试证：在这些条件下，存在 Cantor 三分集 K 到 X 上的一个满射连续映射。

¶ 12) 试由习题 11 推出：

\begin{enumerate}
\def\labelenumi{\alph{enumi})}
\item
若 A 是 R 的一个紧子集，则存在 Cantor 三分集 K 到 A 上的一个连续映射{[}取 A(s) 为 A 与适当选取的区间的交{]}。
\item
若此外 A 还是完备且全不连通的，则它同胚于 K{[}同样的方法，只需这样安排：若 \(s_{1}\) 与 \(s_{2}\) 是两个有相同项数（各项为 o 或 1）的不同序列，则
\end{enumerate}

\[
\mathbf {A} (s _ {1}) \cap \mathbf {A} (s _ {2}) = \varnothing_ {1} ].
\]

¶ 13) 设 X 是一个可数的豪斯多夫空间。假设对每个各项为 o 或 1 的有限序列 s，存在 X 的一个非空子集 B(s)，满足下列条件：

\begin{enumerate}
\def\labelenumi{(\roman{enumi})}
\item
若 \(s\) 是 \(n\) 项（等于 \(\mathbf{o}\) 或 \(\mathbf{i}\) ）的有限序列，且 \(s', s''\) 是 \(n + \mathbf{i}\) 项（等于 \(\mathbf{o}\) 或 \(\mathbf{i}\) ）而前 \(n\) 项与 \(s\) 的相同的序列，则 \(\mathbf{B}(s') \cup \mathbf{B}(s'') = \mathbf{B}(s)\) 且 \(\mathbf{B}(s') \cap \mathbf{B}(s'') = \emptyset\) ；又若 \(s_0\) 是空序列，则 \(\mathbf{B}(s_0) = \mathbf{X}\) 。
\item
对每个 \(x \in \mathbf{X}\) ，若 \(s_n\) 是 \(n\) 项（等于 0 或 1）的（唯一的）序列，它满足 \(x \in \mathrm{B}(s_n)\) ，则由诸 \(\mathrm{B}(s_n)\) 形成的滤子基收敛于 \(x\) （在 \(\mathbf{X}\) 中）。
\end{enumerate}

试证：在这些条件下，E 同胚于有理直线 Q。{[}首先证明 X 同胚于 Cantor 三分集 K 的一个可数子集，它在 K 中稠密且不含 K 的任何邻接区间的端点。为此，注意假设使每个 \(x \in X\) 对应一个各项等于 0 或 1 的无穷序列 \((u_{n}(x))\) ；利用 X 可数这一事实，证明双射 \(s \to \mathbf{B}(s)\) 可以这样修改，使得对任何 \(x \in E\) ，诸 \(u_{n}(x)\) 都不构成最终常数的序列；最后用 §2 的习题 10{]}。

由此推出：\(\mathbf{R}\) 的每个没有孤立点的可数子空间同胚于 \(\mathbf{Q}\) 。

\begin{enumerate}
\def\labelenumi{\arabic{enumi})}
\setcounter{enumi}{13}
\tightlist
\item
试证：\(\mathbf{R}\) 的每个闭子集或者是可数的，或者具有连续统的势{[}用上面的习题 12 b) 以及第一章，§ 9 的习题 16{]}。
\end{enumerate}


\begin{enumerate}
\def\labelenumi{\arabic{enumi})}
\setcounter{enumi}{14}
\item
试证：\(\mathbf{R}\) 的开子集所成之集与 \(\mathbf{R}\) 的紧、全不连通、完备子集所成之集都具有连续统的势。
\item
  \begin{enumerate}
  \def\labelenumii{\alph{enumii})}
  \tightlist
  \item
设 A 是 R 的一个可数闭子集，f 是定义在 R 上的一个连续实值函数。试证：若 f 在邻接于 A 的每个区间中是常数，则 f 是常数（用 Bolzano 定理）。
  \end{enumerate}
\end{enumerate}

\begin{enumerate}
\def\labelenumi{\alph{enumi})}
\setcounter{enumi}{1}
\tightlist
\item
若 \(f\) 是习题 9 中定义的 Cantor 三分集 \(\mathbf{K}\) 到 \([0, 1]\) 上的那个连续映射，试证：\(f\) 可以延拓为 \(\mathbf{R}\) 上的一个连续函数，它在邻接于 \(\mathbf{K}\) 的每个区间上是常数。
\end{enumerate}

\begin{enumerate}
\def\labelenumi{\arabic{enumi})}
\setcounter{enumi}{16}
\tightlist
\item
设 X 是有一个可数稠密子集的拓扑空间。试证：X 上所有实值连续函数所成之集具有连续统的势。
\end{enumerate}
% semantic-fragments: legacy-input-seam
\relax{}
